\documentclass[11pt,reqno]{amsart}

\usepackage[T1]{fontenc}
\usepackage{lmodern}
\usepackage{microtype}
\usepackage{mathtools,amssymb,amsfonts}
\usepackage{mathrsfs}
\usepackage{booktabs}
\usepackage{enumitem}
\usepackage[margin=1.02in]{geometry}
\usepackage{xcolor}
\definecolor{linkblue}{HTML}{244E73}
\definecolor{citegreen}{HTML}{315B45}
\usepackage[
  colorlinks=true,
  linkcolor=linkblue,
  citecolor=citegreen,
  urlcolor=linkblue,
  pdfauthor={ChatGPT},
  pdftitle={Counterexamples to Makeev's regular-simplex conjecture in dimensions four and five}
]{hyperref}
\usepackage[nameinlink,noabbrev]{cleveref}
\usepackage{listings}

\allowdisplaybreaks
\setlist[enumerate]{leftmargin=2.2em,itemsep=0.25em,topsep=0.4em}
\lstdefinestyle{certificate}{
  basicstyle=\ttfamily\scriptsize,
  columns=fullflexible,
  keepspaces=true,
  breaklines=true,
  breakatwhitespace=false,
  showstringspaces=false,
  frame=single,
  rulecolor=\color{black!25},
  xleftmargin=0.5em,
  xrightmargin=0.5em
}

\newtheorem{theorem}{Theorem}[section]
\newtheorem{proposition}[theorem]{Proposition}
\newtheorem{lemma}[theorem]{Lemma}
\theoremstyle{remark}
\newtheorem{remark}[theorem]{Remark}

\crefname{theorem}{theorem}{theorems}
\Crefname{theorem}{Theorem}{Theorems}
\crefname{lemma}{lemma}{lemmas}
\Crefname{lemma}{Lemma}{Lemmas}
\crefname{proposition}{proposition}{propositions}
\Crefname{proposition}{Proposition}{Propositions}
\crefname{remark}{remark}{remarks}
\Crefname{remark}{Remark}{Remarks}

\newcommand{\R}{\mathbb R}
\newcommand{\C}{\mathbb C}
\newcommand{\Q}{\mathbb Q}
\newcommand{\one}{\mathbf 1}
\newcommand{\SO}{\mathrm{SO}}
\newcommand{\eps}{\varepsilon}
\newcommand{\ip}[2]{\left\langle #1,#2\right\rangle}
\newcommand{\norm}[1]{\left\lVert #1\right\rVert}
\newcommand{\abs}[1]{\left\lvert #1\right\rvert}

\title[Counterexamples in dimensions four and five]
{Counterexamples to Makeev's regular-simplex conjecture\\
in dimensions four and five}
\author{ChatGPT}
\thanks{This manuscript and its accompanying certificate generator are
AI-generated.  The dimension-five Magma certificate was independently run by
the human verifier.}
\date{August 2026}

\begin{document}

\begin{abstract}
For $d=4$ and $d=5$ we give an explicit smooth odd function on $S^{d-1}$
which does not agree with any linear functional on any rotated normalized
$A_d$ root system.  In the decomposition
$\R^d=\C^2\oplus\R^{d-4}$, for every sufficiently small nonzero real $\eps$
one may take
\[
 f_\eps(z_1,z_2,x)
 =\operatorname{Re}(z_1^2\overline z_2)
 +\eps\abs{z_1}^2\operatorname{Im}(z_1^2\overline z_2).
\]
The two dimensions share the exact-edge-label reduction, a cubic identity
forced by a circle symmetry, and the compactness argument.  Algebraic
transversality is proved directly in dimension four and by an exact rational
calculation in dimension five.  The latter uses a symmetric multiplication
tensor: after linear elimination it is a unit-ideal test involving 55
quadratics and one cubic in 21 variables.  The accompanying standard-library
Python generator implements the tensor test for every fixed $d\ge4$, emits
the dimension-four and dimension-five Magma certificates, and reproduces the
successfully verified dimension-five file byte for byte.  The complete
direct dimension-four calculation and the generator source are included as
appendices.
\end{abstract}

\maketitle

\section{Introduction}

Kuperberg showed that a constant-width formulation of Makeev's
universal-cover conjecture is equivalent to asking whether every continuous
odd function on $S^{d-1}$ agrees with a linear function on a rotated $A_d$
root system.  He proved the assertion for $d=3$ and showed that the homological
obstruction used there vanishes in higher dimensions \cite{Kuperberg1999}.
The following theorem gives counterexamples in the next two dimensions.

Put $N=d+1$ and
\[
 V_d=\one^\perp\subset\R^N,\qquad
 q_i=e_i-\frac1N\one,\qquad
 u_{ij}=\frac{q_i-q_j}{\sqrt2}
 \quad(0\le i<j\le d).
\]
Thus the $u_{ij}$ are the normalized roots of type $A_d$.

\begin{theorem}[Counterexamples in dimensions four and five]
\label{thm:counterexamples}
For each $d\in\{4,5\}$ there exists $\eps_0(d)>0$ such that, whenever
$0<\abs\eps<\eps_0(d)$, the smooth odd function
\begin{equation}\label{eq:counterexample-function}
 f_\eps(z_1,z_2,x)
 =\operatorname{Re}(z_1^2\overline z_2)
 +\eps\abs{z_1}^2\operatorname{Im}(z_1^2\overline z_2),
 \qquad (z_1,z_2,x)\in S^{d-1}\subset\C^2\oplus\R^{d-4},
\end{equation}
has no isometry $g:V_d\to\R^d$ and $b\in\R^d$ satisfying
\[
 f_\eps(gu_{ij})=b\cdot gu_{ij}
 \qquad(0\le i<j\le d).
\]
Consequently Makeev's conjecture fails in dimensions four and five.
\end{theorem}

A transposition of two simplex vertices induces an orientation-reversing
symmetry of the root system.  Hence the formulations with $\mathrm O(d)$ and
$\SO(d)$ are equivalent.  We freely use either one below.

The proof has one dimension-dependent step.  We first give the common
geometric and variational reductions.  We then encode the remaining
algebraic assertion by a multiplication tensor.  In dimension five the
resulting rational ideal was proved to be the unit ideal by Magma.  In
dimension four the complete direct proof from the original argument is
included in \cref{app:d4-transversality}; the generated tensor certificate
provides an independent computational audit.

\section{Regular-simplex labels}\label{sec:labels}

The simplex vectors satisfy
\begin{equation}\label{eq:simplex-frame}
 \sum_{i=0}^{d}q_i=0,
 \qquad
 \ip{q_i}{q_j}=\delta_{ij}-\frac1N,
 \qquad
 \sum_{i=0}^{d}q_iq_i^T=I_{V_d}.
\end{equation}
In particular $\norm{q_i-q_j}^2=2$, so every $u_{ij}$ lies on the unit
sphere.

For an odd function $F:S^{d-1}\to\R$ and an isometry
$g:V_d\to\R^d$, define the skew matrix
\begin{equation}\label{eq:label-matrix}
 Y_F(g)_{ij}=F(gu_{ij}),
 \qquad Y_F(g)_{ji}=-Y_F(g)_{ij},
\end{equation}
and put
\[
 \Pi=I-\frac1N\one\one^T.
\]

\begin{lemma}[Exact edge labels]\label{lem:exact-labels}
For a real skew $N\times N$ matrix $Y$, the following are equivalent:
\begin{enumerate}
\item $Y_{ij}=c\cdot u_{ij}$ for some $c\in V_d$;
\item $Y_{ij}=\rho_i-\rho_j$ for some $\rho_0,\ldots,\rho_d\in\R$;
\item $\Pi Y\Pi=0$.
\end{enumerate}
\end{lemma}

\begin{proof}
If $Y_{ij}=c\cdot u_{ij}$, take
$\rho_i=(c\cdot q_i)/\sqrt2$.  Conversely, subtract the mean of the
$\rho_i$ and set $c=\sqrt2\sum_i\rho_iq_i$.  The Gram identity
\eqref{eq:simplex-frame} gives $c\cdot q_j=\sqrt2\rho_j$.

If $Y_{ij}=\rho_i-\rho_j$, then
$Y=\rho\one^T-\one\rho^T$, so $\Pi Y\Pi=0$.  Conversely, with
$J_0=N^{-1}\one\one^T$, skew-symmetry gives $J_0YJ_0=0$, and
$\Pi Y\Pi=0$ gives $Y=J_0Y+YJ_0$.  Taking
$\rho_i=N^{-1}\sum_jY_{ij}$ proves $Y_{ij}=\rho_i-\rho_j$.
\end{proof}

We call these matrices \emph{exact}.  If $X$ is skew and $X\one=0$, then
\begin{equation}\label{eq:cycle-annihilates-exact}
 \sum_{i<j}X_{ij}(\rho_i-\rho_j)=0.
\end{equation}
This is the cut/cycle decomposition of complete-graph edge flows
\cite{JiangLimYaoYe2011}.  By \cref{lem:exact-labels}, the equality in
\cref{thm:counterexamples} occurs precisely when $Y_{f_\eps}(g)$ is exact.

\section{The common cubic annihilator}\label{sec:cubic-annihilator}

On $\C^2\oplus\R^{d-4}$ put
\begin{equation}\label{eq:WQRH}
 W(z)=z_1^2\overline z_2,
 \qquad Q=\operatorname{Re}W,
 \qquad R=\operatorname{Im}W,
 \qquad H=\abs{z_1}^2R.
\end{equation}
Consider the circle action and its infinitesimal generator
\begin{equation}\label{eq:circle-action}
 h_\theta(z_1,z_2,x)=(e^{i\theta}z_1,e^{2i\theta}z_2,x),
 \qquad J(z_1,z_2,x)=(iz_1,2iz_2,0).
\end{equation}
Both $Q$ and $H$ are invariant under this action.

Write $p_i=gq_i$ and define
\begin{equation}\label{eq:X-Lambda}
 X(g)_{ij}=\ip{p_i}{Jp_j}_{\R},
 \qquad
 \Lambda_g(Y)=\sum_{i<j}X(g)_{ij}Y_{ij}.
\end{equation}
The matrix $X(g)$ is skew and
$\sum_jX(g)_{ij}=\ip{p_i}{J\sum_jp_j}_{\R}=0$.  Thus
\eqref{eq:cycle-annihilates-exact} shows that $\Lambda_g$ annihilates every
exact label matrix.

Let
\[
 P_3(v)=\sum_{i=0}^{d}\ip{q_i}{v}_{\R}^{3}.
\]
For a homogeneous cubic $C$, write $T_C$ for its symmetric trilinear
polarization and use the invariant inner product
\[
 \ip{C}{D}=\sum_{\alpha,\beta,\gamma=1}^{d}
 T_C(e_\alpha,e_\beta,e_\gamma)
 T_D(e_\alpha,e_\beta,e_\gamma).
\]
For $K\in\mathfrak{so}(d)$, let
$(K\cdot C)(x)=\left.\frac{d}{dt}\right|_{t=0}C(e^{-tK}x)$.

\begin{lemma}[Cubic edge identity]\label{lem:cubic-edge-identity}
For every homogeneous cubic $C$ and $K\in\mathfrak{so}(d)$,
\begin{equation}\label{eq:cubic-edge-identity}
 \ip{C}{K\cdot P_3}
 =2\sqrt2\sum_{i<j}\ip{q_i}{Kq_j}_{\R}\,C(u_{ij}).
\end{equation}
\end{lemma}

\begin{proof}
The Parseval identity in \eqref{eq:simplex-frame} gives
\[
 \ip{C}{K\cdot P_3}=3\sum_iT_C(q_i,q_i,Kq_i).
\]
Since $C(u_{ij})=C(q_i-q_j)/(2\sqrt2)$, the right side of
\eqref{eq:cubic-edge-identity} is
\[
 \sum_{i<j}\ip{q_i}{Kq_j}_{\R}C(q_i-q_j).
\]
The pure terms in its cubic expansion vanish because the row sums of
$(\ip{q_i}{Kq_j})_{ij}$ vanish.  The mixed terms combine to
\[
 -3\sum_{i,j}\ip{q_i}{Kq_j}_{\R}T_C(q_i,q_i,q_j).
\]
Finally $Kq_i=-\sum_j\ip{q_i}{Kq_j}_{\R}q_j$, again by Parseval, and the
last display is $3\sum_iT_C(q_i,q_i,Kq_i)$.
\end{proof}

\begin{lemma}[Global cubic annihilation]\label{lem:global-annihilation}
For every $d\ge4$ and every isometry $g:V_d\to\R^d$,
\begin{equation}\label{eq:global-annihilation}
 \Lambda_g(Y_Q(g))=0.
\end{equation}
\end{lemma}

\begin{proof}
Take $K=g^{-1}Jg$ and $C=Q\circ g$ in
\cref{lem:cubic-edge-identity}.  The orthogonal action on cubic tensors is
infinitesimally skew-adjoint, while
$K\cdot(Q\circ g)=(J\cdot Q)\circ g=0$ by circle invariance.  Hence
$\ip{Q\circ g}{K\cdot P_3}=0$, and the edge identity gives the claim.
\end{proof}

The remaining statement is algebraic.

\begin{proposition}[Algebraic transversality]\label{prop:transversality}
For $d=4$ and $d=5$,
\begin{equation}\label{eq:transversality}
 Y_Q(g)\text{ exact}
 \quad\Longrightarrow\quad
 \Lambda_g(Y_H(g))\ne0.
\end{equation}
\end{proposition}

The direct proof for $d=4$ is \cref{app:d4-transversality}.  The common
algebraic encoding and the exact $d=5$ computation occupy the next two
sections.

\section{The cubic exactness locus}\label{sec:exactness-locus}

Write
\[
 p_i=(a_i,b_i,c_i^{(4)},\ldots,c_i^{(d-1)})
 \in\C^2\oplus\R^{d-4}.
\]
Thus $a,b\in\C^N$, while each $c^{(r)}\in\R^N$.  The real frame identities
say that
\[
 \one,a,b,\overline a,\overline b,c^{(4)},\ldots,c^{(d-1)}
\]
are mutually orthogonal in $\C^N$.  Their squared norms are respectively
$N,2,2,2,2,1,\ldots,1$.

At $gu_{ij}=(p_i-p_j)/\sqrt2$, the complex cubic has, apart from the common
factor $1/(2\sqrt2)$, the edge label
\[
 \widetilde Y_{ij}=(a_i-a_j)^2(\overline b_i-\overline b_j).
\]
For $(u\wedge v)_{ij}=u_iv_j-v_iu_j$, expansion modulo exact matrices gives
\begin{equation}\label{eq:omega-universal}
 \omega:=\Pi\widetilde Y\Pi
 =-a^2\wedge\overline b+2a\wedge(a\overline b).
\end{equation}
The $Q$-labels are exact precisely when $\operatorname{Re}\omega=0$.

Define
\begin{align*}
 s&=\frac12\sum_i\abs{a_i}^2a_i,&
 m&=\sum_i a_i^2\overline b_i,&
 n&=\sum_i a_i^3,&
 \ell&=\frac12\sum_i a_i^2b_i,\\
 \lambda&=\frac12\sum_i\abs{a_i}^2\overline b_i,&
 u&=\frac12\sum_i a_i\overline b_i^{2},&
 \delta&=\frac12\sum_i a_i\abs{b_i}^2,\\
 r_r&=\sum_i c_i^{(r)}a_i^2,&
 w_r&=\sum_i c_i^{(r)}a_i\overline b_i
 &&(4\le r<d).
\end{align*}
Expansion in the displayed orthogonal basis gives
\begin{align}
 a^2={}&sa+\frac m2b+\frac n2\overline a+\ell\overline b
       +\sum_{r=4}^{d-1}r_r c^{(r)},\label{eq:a2-full}\\
 a\overline b={}&\lambda a+ub+\frac m2\overline a+\delta\overline b
       +\sum_{r=4}^{d-1}w_r c^{(r)}.\label{eq:ab-full}
\end{align}
Substitution in \eqref{eq:omega-universal}, followed by comparison with the
complex conjugate, proves the next criterion.

\begin{proposition}[Dimension-independent exactness criterion]
\label{prop:universal-exactness}
The $Q$-labels are exact if and only if
\begin{equation}\label{eq:exactness-moments}
 m\in\R,\qquad 4u=\overline n,\qquad
 r_r=w_r=0\ (4\le r<d),\qquad 2\delta=s.
\end{equation}
Equivalently, $m$ is real and the two coordinatewise vector identities
\begin{align}
 a^2&=sa+\frac m2b+\frac n2\overline a+\ell\overline b,
 \label{eq:universal-a2}\\
 a\overline b&=\lambda a+\frac{\overline n}{4}b
 +\frac m2\overline a+\frac s2\overline b
 \label{eq:universal-ab}
\end{align}
hold.
\end{proposition}

\begin{proof}
Substituting \eqref{eq:a2-full}--\eqref{eq:ab-full} into
\eqref{eq:omega-universal} yields
\begin{align*}
 \omega={}&2u\,a\wedge b+2\sum_rw_r\,a\wedge c^{(r)}
 +m\,a\wedge\overline a+(2\delta-s)a\wedge\overline b\\
 &-\frac m2b\wedge\overline b-\sum_rr_r c^{(r)}\wedge\overline b
 -\frac n2\overline a\wedge\overline b.
\end{align*}
The wedges of the members of the orthogonal basis are linearly independent.
Thus $\omega+\overline\omega=0$ is equivalent to
\eqref{eq:exactness-moments}.  Substitution in
\eqref{eq:a2-full}--\eqref{eq:ab-full} then gives
\eqref{eq:universal-a2}--\eqref{eq:universal-ab}, and the converse follows by
reversing the calculation.
\end{proof}

There are $4d-11$ real linear conditions in
\eqref{eq:exactness-moments}: five common conditions and four for every extra
real coordinate.  This linearity becomes particularly useful in tensor
coordinates.

\section{The multiplication-tensor certificate}\label{sec:tensor}

Write $a=x+iy$ and $b=p+iq$, and set
\[
 e_*=N^{-1/2}\one,\qquad
 v_0=x,\quad v_1=y,\quad v_2=p,\quad v_3=q,\quad
 v_r=c^{(r)}\ (4\le r<d).
\]
Then $e_*,v_0,\ldots,v_{d-1}$ is an orthonormal basis of $\R^N$.  Define the
fully symmetric tensor
\begin{equation}\label{eq:tensor-definition}
 T_{ijk}=\sum_{\nu=0}^{d}v_i(\nu)v_j(\nu)v_k(\nu),
 \qquad 0\le i,j,k<d.
\end{equation}
Coordinatewise multiplication by $v_i$, in this basis, has the real
symmetric matrix
\begin{equation}\label{eq:multiplication-matrix}
 M_i=
 \begin{pmatrix}
 0&N^{-1/2}(\delta_{ij})_{j=0}^{d-1}\\
 N^{-1/2}(\delta_{ik})_{k=0}^{d-1}&(T_{ijk})_{j,k=0}^{d-1}
 \end{pmatrix}.
\end{equation}
Coordinatewise multiplication gives
\begin{equation}\label{eq:trace-commute}
 \sum_{j=0}^{d-1}T_{ijj}=0,
 \qquad [M_i,M_j]=0\quad(0\le i<j<d).
\end{equation}

\begin{lemma}[Exact tensor model]\label{lem:tensor-converse}
A real symmetric tensor $T$ arises from an orthonormal simplex frame by
\eqref{eq:tensor-definition} if and only if it satisfies
\eqref{eq:trace-commute}.
\end{lemma}

\begin{proof}
Necessity was just observed.  Conversely, use $T$ to define a commutative
algebra with unit $\sqrt N e_*$ and multiplication table
\[
 v_iv_j=\frac{\delta_{ij}}{\sqrt N}e_*+\sum_{k=0}^{d-1}T_{ijk}v_k.
\]
The matrices $M_i$ are its left-multiplication operators.  Their commutation,
together with commutativity of the product, implies associativity.  Hence
\begin{equation}\label{eq:multiplication-identity}
 M_iM_j=\frac{\delta_{ij}}N I_N+\sum_{k=0}^{d-1}T_{ijk}M_k.
\end{equation}
The trace equations give $\operatorname{tr}M_k=0$, so taking traces in
\eqref{eq:multiplication-identity} yields
$\operatorname{tr}(M_iM_j)=\delta_{ij}$.

The commuting real symmetric matrices are simultaneously orthogonally
diagonalizable.  Their joint eigenvalue vectors are therefore an orthonormal
basis of $\one^\perp\subset\R^N$ (zero trace gives orthogonality to $\one$).
Finally,
\[
 \operatorname{tr}(M_iM_jM_k)=T_{ijk},
\]
by \eqref{eq:multiplication-identity}; in the common eigenbasis this trace is
exactly the right side of \eqref{eq:tensor-definition}.  Thus no real frames
are lost and no spurious real frames are introduced.
\end{proof}

Only the lower-right block of each commutator is nontrivial.  Its independent
entries are
\begin{equation}\label{eq:commutator-polynomials}
 \sum_{m=0}^{d-1}
 \bigl(T_{ikm}T_{jm\ell}-T_{jkm}T_{im\ell}\bigr)
 +\frac{\delta_{ki}\delta_{j\ell}-\delta_{kj}\delta_{i\ell}}N=0,
\end{equation}
where $i<j$ and $k<\ell$.

In zero-based indices $0,1,2,3=x,y,p,q$, the exactness equations of
\cref{prop:universal-exactness} are the following rational linear equations:
\begin{align}
 2T_{012}-T_{003}+T_{113}&=0,\notag\\
 2T_{022}-2T_{033}+4T_{123}-T_{000}+3T_{011}&=0,\notag\\
 2T_{122}-2T_{133}-4T_{023}+3T_{001}-T_{111}&=0,\notag\\
 T_{00r}-T_{11r}=2T_{01r}=T_{02r}+T_{13r}
 =T_{12r}-T_{03r}&=0\quad(4\le r<d),\notag\\
 2T_{022}+2T_{033}-T_{000}-T_{011}&=0,\notag\\
 2T_{122}+2T_{133}-T_{001}-T_{111}&=0.
\label{eq:tensor-exactness-linear}
\end{align}
Together with the $d$ trace equations, these are $5d-11$ independent linear
conditions in the $\binom{d+2}{3}$ tensor coordinates.  Exact rational row
reduction therefore leaves
\begin{equation}\label{eq:number-free}
 r_d=\binom{d+2}{3}-(5d-11)
\end{equation}
free variables.

It remains to express the obstruction.  In addition to $s,m,n,\ell,\lambda$
from \cref{sec:exactness-locus}, put
\[
 v=\sum_i a_i b_i^2,
 \qquad k=\sum_i b_i^3,
 \qquad \rho=\sum_i\abs{b_i}^2b_i.
\]
All eight moments are linear in $T$:
\begin{align*}
 s={}&\tfrac12(T_{000}+T_{011})
       +\tfrac i2(T_{001}+T_{111}),\\
 m={}&T_{002}-T_{112}+2T_{013},\\
 n={}&(T_{000}-3T_{011})+i(3T_{001}-T_{111}),\\
 \ell={}&\tfrac12(T_{002}-T_{112}-2T_{013})
       +\tfrac i2(T_{003}-T_{113}+2T_{012}),\\
 \lambda={}&\tfrac12(T_{002}+T_{112})
       -\tfrac i2(T_{003}+T_{113}),\\
 v={}&(T_{022}-T_{033}-2T_{123})
       +i(T_{122}-T_{133}+2T_{023}),\\
 k={}&(T_{222}-3T_{233})+i(3T_{223}-T_{333}),\\
 \rho={}&(T_{222}+T_{233})+i(T_{223}+T_{333}).
\end{align*}
On the linear exactness locus, $m$ is real.  Define
\begin{align}
 \mathscr E={}&-\tfrac32m^3+\tfrac34\bar nmn
 -12\bar\rho\ell\lambda-16\bar\lambda\lambda m
 -2\bar n\bar s\ell-6\rho\lambda m
 -8\bar s\lambda n-8\lambda s^2-\bar v m v\notag\\
 &+12\bar k\bar\lambda\ell+14\bar\ell\ell m
 +16\ell\lambda^2+4\bar\lambda\bar v n
 +5\bar\ell ns+6\bar\lambda\bar\rho m+6\bar v\ell s,
 \qquad E=\operatorname{Re}\mathscr E.
\label{eq:obstruction-E}
\end{align}

\begin{lemma}[Tensor obstruction identity]\label{lem:Escale}
On the cubic exactness locus, in every dimension $d\ge4$,
\begin{equation}\label{eq:Escale}
 E=8\sqrt2\,\Lambda_g(Y_H(g)).
\end{equation}
\end{lemma}

\begin{proof}
At an edge, the degree-five factor in $H$ is
\[
 \frac1{4\sqrt2}(a_i-a_j)^3
 (\overline a_i-\overline a_j)(\overline b_i-\overline b_j),
\]
and
\[
 X(g)=\frac i2(\overline a\wedge a+2\overline b\wedge b).
\]
Expand the edge factor modulo exact matrices, pair its wedges with $X(g)$
using
\[
 \sum_{i<j}(u\wedge v)_{ij}(x\wedge y)_{ij}
 =(u\cdot x)(v\cdot y)-(u\cdot y)(v\cdot x),
\]
and reduce every coordinatewise product by
\eqref{eq:universal-a2}--\eqref{eq:universal-ab}.  Collecting the real part
gives exactly \eqref{eq:obstruction-E}.  This calculation is independent of
the number of extra real frame vectors because their coefficients vanish in
the two universal identities.  Keeping the edge normalization gives the
factor $8\sqrt2$ in \eqref{eq:Escale}.  The full uncollected wedge expansion
and its specialization to $d=4$ appear in
\cref{app:d4-transversality}; the generator in the next section
implements the collected rational expression term by term.
\end{proof}

Thus algebraic transversality follows if the ideal generated by
\eqref{eq:commutator-polynomials} and $E$ becomes the unit ideal after the
linear equations have been eliminated.  By \cref{lem:tensor-converse}, this
test loses no real simplex frames.  A unit ideal over $\C$ is stronger than
needed, since it excludes complex as well as real common zeros.

\section{Exact computations in dimensions four and five}
\label{sec:computations}

The accompanying script \texttt{tensor\_polynomials.py} uses only the Python
standard library.  For a selected $d\ge4$ it:
\begin{enumerate}
\item enumerates the symmetric coordinates $T_{ijk}$;
\item performs exact rational row reduction on the $5d-11$ linear equations;
\item substitutes the resulting parametrization into
      \eqref{eq:commutator-polynomials} and \eqref{eq:obstruction-E};
\item clears only nonzero rational scalar factors, deletes zero and duplicate
      commutators, and emits a self-contained Magma program over $\Q$.
\end{enumerate}
The resulting systems are summarized below.

\begin{center}
\begin{tabular}{@{}c@{\quad}r@{\quad}r@{\quad}r@{\quad}r@{\quad}r@{}}
\toprule
$d$&tensor variables&linear equations&free variables&quadratics&cubics\\
\midrule
4&20&9&11&21&1\\
5&35&14&21&55&1\\
\bottomrule
\end{tabular}
\end{center}

The dimension-four obstruction has 73 monomials after substitution and is
$E/4$ up to the generator's harmless normalization.  The dimension-five
obstruction has 163 monomials and is $9E/4$.  Generate and run both files by
\begin{lstlisting}[style=certificate]
python3 tensor_polynomials.py --dimension 4 \
  --emit-magma makeev_d4_tensor_certificate.m
python3 tensor_polynomials.py --dimension 5 \
  --emit-magma makeev_d5_tensor_certificate.m
magma makeev_d4_tensor_certificate.m
magma makeev_d5_tensor_certificate.m
\end{lstlisting}
Each generated program forms the ideal, enables Magma's modular
Gr\"obner-basis arithmetic, calls \texttt{IsProper}, asserts that the answer
is false, and prints
\begin{verbatim}
Certificate verified: tensor ideal is the unit ideal.
\end{verbatim}
only after that assertion passes.

For $d=5$, the generated file is byte-for-byte identical to the certificate
that was run successfully in approximately 14 seconds.  The exact source
digests for the accompanying files are
\begin{center}
\begin{tabular}{@{}l@{}}
\texttt{fe71e120af7f7c38806701e07042f88bff461c5a19738653d78a0fbe9a7677ff}\\
\quad\texttt{tensor\_polynomials.py}\\[2pt]
\texttt{4ae9030dd0b57c9541c7e6861d9d87a956015e009bd5220c2903eb222018d8d7}\\
\quad\texttt{makeev\_d4\_tensor\_certificate.m}\\[2pt]
\texttt{403d831600c2ac4a6d60ba768506ff031eb1d4817eadde04d80e40e3aa14d2e6}\\
\quad\texttt{makeev\_d5\_tensor\_certificate.m}.
\end{tabular}
\end{center}
Therefore the dimension-five part of
\cref{prop:transversality} is an exact, unconditional consequence of that
unit-ideal computation.  The supplied dimension-four Magma file is an
independent audit of the direct proof in \cref{app:d4-transversality}; the
dimension-four theorem does not depend on an unreported computer run.

The generator accepts every fixed $d\ge4$.  This is a way to formulate
higher-dimensional tests, not a claim that their Gr\"obner-basis computations
will terminate quickly.  The number $r_d$ of free variables grows cubically.
In particular, the older coordinate program
\texttt{makeev\_general\_certificate.m}, which is slow already for $d=4,5$,
is not used anywhere in the proof.

\section{Completion of the proof}
\label{sec:completion}

Fix $d\in\{4,5\}$.  Suppose the theorem failed for arbitrarily small
nonzero $\eps$.  Then there would be $\eps_\nu\to0$, $\eps_\nu\ne0$, and
isometries $g_\nu:V_d\to\R^d$ for which $Y_{Q+\eps_\nu H}(g_\nu)$ is exact.
Equivalently,
\begin{equation}\label{eq:defect-equation}
 \Pi Y_Q(g_\nu)\Pi
 +\eps_\nu\Pi Y_H(g_\nu)\Pi=0.
\end{equation}
After passing to a subsequence, compactness of the orthogonal group gives
$g_\nu\to g$.  Letting $\nu\to\infty$ in
\eqref{eq:defect-equation} shows that $Y_Q(g)$ is exact.

For every $\nu$, the functional $\Lambda_{g_\nu}$ annihilates exact labels;
also $\Lambda_{g_\nu}(Y_Q(g_\nu))=0$ by
\cref{lem:global-annihilation}.  Since $\eps_\nu\ne0$, exactness of the
$Q+\eps_\nu H$ labels therefore gives
\[
 \Lambda_{g_\nu}(Y_H(g_\nu))=0.
\]
Passing to the limit yields $\Lambda_g(Y_H(g))=0$.  This contradicts
\cref{prop:transversality}, because $Y_Q(g)$ is exact.  Hence an
$\eps_0(d)>0$ exists.  The polynomial in
\eqref{eq:counterexample-function} is smooth and odd, completing the proof.

\appendix

\section{Direct proof of algebraic transversality in dimension four}
\label{app:d4-transversality}

This appendix reproduces the complete algebraic appendix of the original
dimension-four argument.  It proves the $d=4$ assertion of
\cref{prop:transversality} without computer assistance.

Fix an isometry $g:V_4\to\R^4\cong\C^2$ for which the $Q$-labels are exact,
write
\[
p_i=gq_i=(a_i,b_i)\in\C^2,
\qquad
a=(a_0,\dots,a_4),\quad b=(b_0,\dots,b_4)\in\C^5,
\]
and use coordinatewise multiplication of vectors in $\C^5$.  The Hermitian
inner product is conjugate-linear in the first variable:
\[
\ip{x}{y}_{\C^5}=\sum_i\overline{x_i}y_i.
\]

\subsection{Frame identities}

The $5\times4$ real matrix whose rows are the $p_i$ has orthonormal columns,
all perpendicular to $\one$.  Therefore
\begin{align}
\sum_i a_i=\sum_i b_i&=0,\label{eq:frame-sums}\\
\sum_i a_i^2=\sum_i b_i^2
=\sum_i a_ib_i=\sum_i a_i\overline b_i&=0,\label{eq:frame-orthog}\\
\sum_i\abs{a_i}^2=\sum_i\abs{b_i}^2&=2,\label{eq:frame-norms}\\
\abs{a_i}^2+\abs{b_i}^2&=\frac45.\label{eq:row-norm}
\end{align}
Consequently
\[
\one,a,b,\overline a,\overline b
\]
are mutually orthogonal in $\C^5$, with squared norms $5,2,2,2,2$.
Every $x\in\C^5$ has the expansion
\begin{equation}\label{eq:basis-expansion}
x=\frac{\sum_i x_i}{5}\one
+\frac12\left(
\ip{a}{x}a+\ip{b}{x}b
+\ip{\overline a}{x}\overline a
+\ip{\overline b}{x}\overline b
\right).
\end{equation}
For $u,v\in\C^5$, write
\[
(u\wedge v)_{ij}=u_iv_j-v_iu_j.
\]

\subsection{Cycle equations for the cubic}

At the root $gu_{ij}=(p_i-p_j)/\sqrt2$,
\[
W(gu_{ij})
=\frac1{2\sqrt2}(a_i-a_j)^2(\overline b_i-\overline b_j).
\]
Ignoring the harmless factor $1/(2\sqrt2)$, let
\[
\widetilde Y_{ij}=(a_i-a_j)^2(\overline b_i-\overline b_j).
\]
Modulo exact matrices, expansion gives
\begin{equation}\label{eq:cubic-cycle-form}
\Pi\widetilde Y\Pi
\equiv -a^2\wedge\overline b+2a\wedge(a\overline b).
\end{equation}
Indeed, the omitted term is
$a_i^2\overline b_i-a_j^2\overline b_j$.

Define
\begin{align}
s&=\frac12\sum_i\abs{a_i}^2a_i,
&\lambda&=\frac12\sum_i\abs{a_i}^2\overline b_i,\label{eq:s-lambda}\\
m&=\sum_i a_i^2\overline b_i,
&n&=\sum_i a_i^3,
&u&=\frac12\sum_i a_i\overline b_i^2,
&\ell&=\frac12\sum_i a_i^2b_i.\label{eq:m-n-u-ell}
\end{align}
Using \eqref{eq:basis-expansion},
\begin{align}
a^2&=sa+\frac m2b+\frac n2\overline a+\ell\overline b,
\label{eq:a2-expansion}\\
a\overline b&=\lambda a+ub+\frac m2\overline a-s\overline b.
\label{eq:abbar-expansion}
\end{align}
Substitution into \eqref{eq:cubic-cycle-form} yields
\begin{equation}\label{eq:omega}
\omega:=\Pi\widetilde Y\Pi
=2u\,a\wedge b
+m\,a\wedge\overline a
-3s\,a\wedge\overline b
-\frac m2b\wedge\overline b
-\frac n2\overline a\wedge\overline b.
\end{equation}
Exactness of the real $Q$-labels is precisely $\operatorname{Re}\omega=0$.
The six wedges of the complex basis $a,b,\overline a,\overline b$ are linearly
independent.  Comparing \eqref{eq:omega} with its conjugate gives
\begin{equation}\label{eq:cubic-exactness-conditions}
s=0,
\qquad m\in\R,
\qquad 4u=\overline n.
\end{equation}

Replacing $g$ by $h_\theta g$ changes neither the values of $Q$ and $H$ nor
the matrix $X(g)$, while
$\lambda\mapsto e^{-2i\theta}\lambda$.  We may therefore assume
\begin{equation}\label{eq:t-gauge}
\lambda=t\ge0.
\end{equation}
In addition define
\begin{equation}\label{eq:v-k}
v=\sum_i a_ib_i^2,
\qquad
k=\sum_i b_i^3.
\end{equation}

\subsection{The coordinatewise multiplication table}

Expanding each product in the orthogonal basis
$\one,a,b,\overline a,\overline b$ by \eqref{eq:basis-expansion}, and using
\eqref{eq:row-norm} and \eqref{eq:cubic-exactness-conditions}, gives
\begin{align}
a\overline a&=\frac25\one+tb+t\overline b,
&b\overline b&=\frac25\one-tb-t\overline b,\notag\\
a^2&=\frac m2b+\frac n2\overline a+\ell\overline b,
&\overline a^2&=\frac m2\overline b+\frac{\overline n}{2}a+\overline\ell b,
\notag\\
a\overline b&=ta+\frac{\overline n}{4}b+\frac m2\overline a,
&\overline a b&=t\overline a+\frac n4\overline b+\frac m2a,
\notag\\
ab&=ta+\ell\overline a+\frac v2\overline b,
&\overline a\,\overline b&=t\overline a+\overline\ell a
 +\frac{\overline v}{2}b,
\notag\\
b^2&=\frac n4a-tb+\frac v2\overline a+\frac k2\overline b,
&\overline b^2&=\frac{\overline n}{4}\overline a-t\overline b
 +\frac{\overline v}{2}a+\frac{\overline k}{2}b.
\label{eq:multiplication-table}
\end{align}
For example, the coefficient of $b$ in $a\overline a$ is
$\frac12\sum_i\overline b_i\abs{a_i}^2=t$, and the other coefficients are
obtained in the same way.

Coordinatewise multiplication is associative.  Comparing coefficients in
\eqref{eq:multiplication-table} gives the following necessary equations:
\begin{align}
3mn&=4tv,\tag{E1}\label{eq:E1}\\
\overline n\,v+16\ell t+8mt&=0,\tag{E2}\label{eq:E2}\\
-2km+24\ell t-n^2&=0,\tag{E3}\label{eq:E3}\\
-kt+\ell m+4t^2&=0,\tag{E4}\label{eq:E4}\\
\overline n\,k+mn+8tv&=0,\tag{E5}\label{eq:E5}\\
20\abs\ell^2+5m^2+5\abs n^2-40t^2&=8,\tag{E6}\label{eq:E6}\\
16\abs\ell^2-4m^2-\abs n^2+4\abs v^2&=0,\tag{E7}\label{eq:E7}\\
20\abs\ell^2+60t^2+5\abs v^2&=8.\tag{E8}\label{eq:E8}
\end{align}
More precisely, (E1)--(E3) are the $a,b,\overline b$ coefficients of
$(a^2)b=a(ab)$; (E4)--(E5) are the $a,b$ coefficients of
$(ab)b=a(b^2)$; and (E6), (E7), (E8) are respectively the indicated
coefficients of
\[
(a^2)\overline a=a(a\overline a),\qquad
(ab)\overline a=a(b\overline a),\qquad
(ab)\overline b=a(b\overline b).
\]

\subsection{Classification of the possible cubic zeros}

We extract the consequences of (E1)--(E8) needed below.

\medskip
\noindent\textbf{Case I: $t=0$.}
Equation (E1) gives $mn=0$.  If $m=0$, then (E3) gives $n=0$, while
(E6)--(E7) become
\[
20\abs\ell^2=8,
\qquad
16\abs\ell^2+4\abs v^2=0,
\]
a contradiction.  Hence $m\ne0$ and $n=0$.  Equations (E4) and (E3) give
$\ell=k=0$, and then
\begin{equation}\label{eq:case-I}
m^2=\frac85,
\qquad
\abs v^2=m^2.
\end{equation}

\medskip
\noindent\textbf{Case II: $t>0$ and $n=0$.}
Equation (E1) gives $v=0$, and (E2) gives
\begin{equation}\label{eq:case-II-ell}
\ell=-\frac m2.
\end{equation}
Here $m\ne0$, since otherwise \eqref{eq:multiplication-table} would give
$a^2=0$.  Equations (E3) and (E4) give
$k=-6t$ and $m^2=20t^2$.  Finally (E8) yields
\begin{equation}\label{eq:case-II}
t^2=\frac1{20},
\qquad
m^2=1.
\end{equation}

\medskip
\noindent\textbf{Case III: $t>0$ and $n\ne0$.}
First $m\ne0$, since otherwise (E1), (E2), and (E3) successively give
$v=0$, $\ell=0$, and $n=0$.  Put
\begin{equation}\label{eq:q-h}
q=\frac mt,
\qquad
h=\frac{\abs n^2}{t^2}>0.
\end{equation}
Equations (E1) and (E2) give
\begin{equation}\label{eq:v-ell-case-III}
v=\frac{3mn}{4t},
\qquad
\ell=-\frac{m(32+3h)}{64}.
\end{equation}
In particular, $\ell$ is real.  Equation (E5) gives
\[
k=-7m\frac n{\overline n}.
\]
Equation (E4) shows that $k$ is real, so
\begin{equation}\label{eq:sigma}
\sigma:=\frac n{\overline n}\in\{1,-1\},
\qquad
k=-7m\sigma,
\qquad
n^2=\sigma\abs n^2.
\end{equation}
After division by appropriate powers of $t$, equations (E4) and (E3) become
\begin{align}
3hq^2+32q^2-448\sigma q-256&=0,\label{eq:case-III-a}\\
-9hq-8\sigma h+112\sigma q^2-96q&=0.\label{eq:case-III-b}
\end{align}
Eliminating $h$ gives
\begin{align}
(q-4)(q+2)(21q^2+42q+16)&=0 &&(\sigma=1),\label{eq:resultant-plus}\\
(q+4)(q-2)(21q^2-42q+16)&=0 &&(\sigma=-1).
\label{eq:resultant-minus}
\end{align}
The roots $q=-2\sigma$ give $h=-64$ and are impossible.  One root of each
quadratic gives
$h=(-74-14\sqrt{105})/3<0$.  The candidates $q=4\sigma$ give $h=32$;
then \eqref{eq:v-ell-case-III} and (E1) imply
\[
\frac{\abs\ell^2}{t^2}=64,
\qquad
\frac{\abs v^2}{t^2}=288,
\]
and the left side of (E7), divided by $t^2$, is $2080$, a contradiction.
The only remaining possibility is
\begin{equation}\label{eq:case-III-qh}
q=-\sigma\left(1-\frac{\sqrt{105}}{21}\right),
\qquad
h=\frac{-74+14\sqrt{105}}3.
\end{equation}
Equation (E6) then gives
\begin{equation}\label{eq:case-III-tm}
t^2=\frac{7(5\sqrt{105}-33)}{1440},
\qquad
m^2=\frac{49\sqrt{105}-477}{1080}.
\end{equation}
These three cases exhaust every actual cubic zero, because every such zero
satisfies the necessary equations (E1)--(E8).

\subsection{Evaluation of the quintic functional}

For the actual root $gu_{ij}=(p_i-p_j)/\sqrt2$,
\begin{equation}\label{eq:quintic-edge}
\abs{z_1}^2W(z)
=\frac1{4\sqrt2}
(a_i-a_j)^3(\overline a_i-\overline a_j)
(\overline b_i-\overline b_j).
\end{equation}
Define
\begin{equation}\label{eq:Sigma}
\Sigma=\sum_{i<j}X(g)_{ij}
(a_i-a_j)^3(\overline a_i-\overline a_j)
(\overline b_i-\overline b_j).
\end{equation}
Then
\begin{equation}\label{eq:Lambda-H-Sigma}
\Lambda_g(Y_H(g))=\frac1{4\sqrt2}\operatorname{Im}\Sigma.
\end{equation}
In wedge notation,
\begin{equation}\label{eq:X-wedge}
X(g)=\frac i2\bigl(\overline a\wedge a+2\overline b\wedge b\bigr).
\end{equation}
Modulo an exact matrix, the degree-five edge expression in
\eqref{eq:Sigma} is
\begin{align}
F\equiv{}&
\overline b\wedge(a^3\overline a)
+\overline a\wedge(a^3\overline b)
-(\overline a\,\overline b)\wedge a^3\notag\\
&+3a\wedge(a^2\overline a\,\overline b)
-3(a\overline b)\wedge(a^2\overline a)\notag\\
&-3(a\overline a)\wedge(a^2\overline b)
+3(a\overline a\,\overline b)\wedge a^2.
\label{eq:F-wedge}
\end{align}
The omitted exact term is
$(a^3\overline a\,\overline b)\wedge\one$, which is annihilated by $X(g)$.
For the bilinear dot product $x\cdot y=\sum_i x_iy_i$,
\begin{equation}\label{eq:wedge-pairing}
\sum_{i<j}(u\wedge v)_{ij}(x\wedge y)_{ij}
=(u\cdot x)(v\cdot y)-(u\cdot y)(v\cdot x).
\end{equation}
Every dot product in \eqref{eq:wedge-pairing} is five times the coefficient of
$\one$ in the corresponding coordinatewise product.  Substituting
\eqref{eq:multiplication-table}, \eqref{eq:X-wedge}, and
\eqref{eq:F-wedge} gives
\begin{equation}\label{eq:Sigma-Braw}
\Sigma=-\frac i{80}\,\mathcal B_{\mathrm{raw}},
\end{equation}
where
\begin{align}
\mathcal B_{\mathrm{raw}}={}&
-960\overline k\,\ell t
+160\abs\ell^2m
+320\ell t^2
-120m^3\notag\\
&-75m\abs n^2
-800mt^2
+100m\abs v^2
+192m
-160nt\overline v.
\label{eq:Braw}
\end{align}
In all three cases above, $\ell$ is real.  The conjugates of (E4) and (E1)
give
\begin{equation}\label{eq:two-substitutions}
\overline k\,t=\ell m+4t^2,
\qquad
nt\overline v=\frac34m\abs n^2.
\end{equation}
Using \eqref{eq:two-substitutions}, followed by (E6) and (E8), reduces
\eqref{eq:Braw} to
\begin{equation}\label{eq:B-simplified}
\mathcal B
=m\bigl(180m^2+105\abs n^2-4400t^2-128\bigr)
-3520\ell t^2.
\end{equation}
The three cases now give
\begin{equation}\label{eq:B-values}
\begin{array}{c@{\qquad}c}
\toprule
\text{case}&\mathcal B\\
\midrule
\mathrm{I}&160m\\
\mathrm{II}&-80m\\
\mathrm{III}&60m(75-7\sqrt{105})\\
\bottomrule
\end{array}
\end{equation}
Each value is nonzero: in every case $m\ne0$, and
$75-7\sqrt{105}>0$.  Since $\mathcal B$ is real,
\eqref{eq:Lambda-H-Sigma} and \eqref{eq:Sigma-Braw} yield
\[
\Lambda_g(Y_H(g))=-\frac{\mathcal B}{320\sqrt2}\ne0.
\]
This proves the dimension-four assertion of
\cref{prop:transversality}.

\begin{thebibliography}{9}

\bibitem{JiangLimYaoYe2011}
X.~Jiang, L.-H.~Lim, Y.~Yao, and Y.~Ye,
\emph{Statistical ranking and combinatorial Hodge theory},
Math. Program. \textbf{127} (2011), 203--244.
\href{https://doi.org/10.1007/s10107-010-0419-x}{doi:10.1007/s10107-010-0419-x}.

\bibitem{Kuperberg1999}
G.~Kuperberg,
\emph{Circumscribing constant-width bodies with polytopes},
New York J. Math. \textbf{5} (1999), 91--100.
\url{https://nyjm.albany.edu/j/1999/5-6p.pdf}.

\end{thebibliography}

\end{document}
